\section{Paradigma Divide and Conquer}
\subsection{Proprietà di sottostruttura}
\subsubsection*{Idea generale}
Il paradigma divide and conquer (divide et impera) è un approccio che consiste nel suddividere un'istanza di grandi dimensioni
in una o più istanze di dimensioni inferiori, più semplici da risolvere, e poi combinare le soluzioni di queste istanze per
ottenere la soluzione dell'istanza originale.

\subsubsection*{Proprietà di sottostruttura}
Dato un problema \(\pi \subseteq \mathcal{I} \times \mathcal{S}\), il paradigma si compone di tre componenti:
\begin{itemize}[topsep=5pt, itemsep=3pt]
	\item \(A_D : \mathcal{I} \to \mathcal{I}*\), funzione di divide che suddivide un'istanza \(i \in \mathcal{I}\) di
	taglia \(|i| > n_0\) in una sequenza di una o più istanze più piccole \(<\!i_1, i_2, \dots i_k\!>\) di taglia
	\(|i_j| < |i|\) per ogni \(j\)
	\item \(A_C : \mathcal{S}* \to \mathcal{S}\), funzione di combine che combina la sequenza di soluzioni
	\(<\!s_1, s_2, \dots s_k\!>\) associate alle istanze \(<\!i_1, i_2, \dots i_k\!>\) nella soluzione \(s\) associata
	all'istanza \(i\)
	\item \(D\&C : \mathcal{I} \to \mathcal{S}\), chiamata ricorsiva che trasforma individualmente ogni istanza	\(i_j\)
	nella corrispettiva soluzione \(s_j\), applicando eventualmente le due funzioni \(A_D\) e \(A_C\) in modo ricorsivo
\end{itemize}

\noindent
Le due funzioni \(A_D\) e \(A_C\) devono soddisfare la seguente proprietà di sottostruttura:
\[\forall i \in \mathcal{I}, |i| > n_0 \text{ vale } \left[ \; D(i) = \; <\!i_1, \dots, i_k\!> \; \wedge \; \forall j \; (i_j, s_j) \in \pi \; \wedge \; C(<\!s_1, \dots, s_k\!>) = s \implies (i, s) \in \pi \; \right]\]

\subsubsection*{Chiamata ricorsiva (conquer) e definizione del caso base}
La chiamata ricorsiva \(D\&C\) è definita come segue:
\begin{itemize}
	\item se \(|i| > n_0\) allora \(D\&C(i) = A_C(A_D(i))\), ovvero riapplica il modello di divide and conquer
	\item se \(|i| \leq n_0\) allora utilizza una soluzione diretta ed immediata per risolvere \(i\)
\end{itemize}
In generale il parametro \(n_0\) è scelto in base al problema. Può essere il minimo valore per cui è possibile applicare
il modello di divide and conquer, oppure un valore sufficientemente grande sotto di cui si conosce un algoritmo più
efficiente per risolvere il problema.

\subsubsection*{Pseudocodice}
\begin{algo}{D\&C}{i}{\text{istanza \(i \in \mathcal{I}\)}}{\text{soluzione \(s \in \mathcal{S}\)}}
\If{\(|i| \leq n_0\)} \Comment{caso base \(\qquad\qquad\qquad\qquad\qquad\qquad\qquad\qquad\qquad\qquad\qquad\;\)}
	\State \(s \gets A_\text{diretto}(i)\) \Comment{soluzione diretta ed immediata \(\quad\qquad\qquad\qquad\qquad\quad\,\)}
	\State \textbf{return} \(s\)
\Else \Comment{passo ricorsivo \(\qquad\qquad\qquad\qquad\qquad\qquad\qquad\qquad\qquad\qquad\)}
	\State \(<\!i_1, i_2, \dots i_k\!> \; \gets A_D(i)\) \Comment{divide \(\quad\qquad\qquad\qquad\qquad\qquad\qquad\qquad\qquad\qquad\qquad\)}
	\For{\(j \gets 1\)}{\(k\)} \Comment{conquer \(\quad\qquad\qquad\qquad\qquad\qquad\qquad\qquad\qquad\qquad\quad\:\)}
		\State \(s_j \gets D\&C(i_j)\)
	\EndFor
	\State \(s \gets A_C(<\!s_1, s_2, \dots s_k\!>)\) \Comment{combine \(\quad\qquad\qquad\qquad\qquad\qquad\qquad\qquad\qquad\qquad\quad\)}
	\State \textbf{return} \(s\)
\EndIf
\end{algo}

\newpage

\subsection{Analisi della correttezza}
\subsubsection*{Proprietà di sottostruttura in funzione di un parametro intero}
Per dimostrare la correttezza di un algoritmo divide and conquer, si utilizza la proprietà di sottostruttura e il principio
di induzione in forma forte. Per prima cosa si esprime la proprietà di sottostruttura in funzione di un parametro intero,
in modo da renderla più adatta alla dimostrazione induttiva, ovvero:
\[\forall i \in \mathcal{I} \; \underbrace{\left[ \; (i, \; D\&C(i)) \in \pi \; \right]}_{P(i)} \quad = \quad \forall n \in \mathbb{N} \; \underbrace{\left[ \forall i \in \mathcal{I} \text{ con } |i| = n  \text{ vale } (i, D\&C(i)) \in \pi \; \right]}_{P(n)}\]

\subsubsection*{Dimostrazione con il principio di induzione in forma forte}
Si dimostra, infine, che \(P(n)\) è vera per ogni \(n \in \mathbb{N}\) tramite il principio di induzione in forma forte:
\begin{itemize}
	\item si dimostra che \(P(n)\) è vera per ogni \(n \leq n_0\) (caso base)
	\item si dimostra che se \(P(n)\) è vera per tutti i valori \(n \leq m\) allora anche \(P(m + 1)\) è vera (passo induttivo)
\end{itemize}

\subsection{Albero della ricorsione}
\subsubsection*{Albero della ricorsione}
L'albero della ricorsione è una rappresentazione grafica che rappresenta le chiamate ricorsive di un algoritmo divide and
conquer con le seguenti caratteristiche:
\begin{itemize}
	\item ogni nodo rappresenta una chiamata ricorsiva \(D\&C(i)\) per un'istanza \(i \in \mathcal{I}\)
	\item ogni nodo ha come figli i nodi che rappresentano le chiamate ricorsive \(D\&C(i_j)\) per ogni istanza
	\(i_j\) ottenuta dalla procedura di divide \(A_D(i)\)
	\item la taglia dell'istanza dei nodi figli è inferiore alla taglia dell'istanza del nodo padre, in quanto la taglia
	dell'istanza diminuisce ad ogni chiamata ricorsiva
	\item ogni foglia rappresenta una chiamata ricorsiva \(D\&C(i)\) per un'istanza \(i \in \mathcal{I}\) di taglia
	\(|i| \leq n_0\), ovvero un caso base
\end{itemize}

\subsubsection*{Esplorazione dell'albero della ricorsione}
L'esplorazione dell'albero della ricorsione avviene secondo la visita in-order, ovvero prima si visita il figlio sinistro,
poi il nodo padre, poi il figlio destro. In questo modo, dato il nodo dell'albero che si sta attualmente visitando, si è
certi che tutti i nodi a sinistra di esso sono stati visitati, e i nodi a destra di esso sono ancora da visitare.

La costruzione della soluzione avviene secondo la visita post-order, ovvero prima si risolvono le istanze dei figli, poi
viene costruita la soluzione per il nodo padre, secondo l'algoritmo DFS (Depth First Search).

Lo spazio occupato è proporzionale allo spazio richiesto da ciascuna chiamata ricorsiva e dal numero di chiamate ricorsive
attive contemporaneamente, ovvero la profondità dell'albero della ricorsione.

\subsubsection*{Cenni sulla ricorsione nei primi linguaggi di programmazione}
La memoria a disposizione sui primi computer era molto limitata e non era sufficiente per implementare uno stack che isolasse
le variabili locali delle varie chiamate ricorsive attive. I primi linguaggi di programmazione, infatti, salvavano le variabili
locali in posizioni fisse di memoria, comuni ad ogni chiamata ricorsiva. Per questo motivo, all'inizio la ricorsione non era
supportata ed il suo utilizzo era fortemente scoraggiato. Il primo linguaggio di programmazione a supportare nativamente la
ricorsione è stato il linguaggio LISP, sviluppato da John McCarthy nel 1958.

Ad oggi, la ricorsione è supportata nativamente da tutti i principali linguaggi di programmazione ed esistono ottimizzazioni
a livello hardware (istruzioni macchina per la gestione dello stack) per renderla efficiente e non sprecare memoria.

\subsection{Analisi della complessità e formulone}
\subsubsection*{Equazione alle ricorrenze}
L'equazione alle ricorrenze è un'equazione che esprime la complessità temporale di un algoritmo divide and conquer
in funzione della complessità temporale delle chiamate ricorsive e delle operazioni di divide e combine. Di seguito
è riportata la forma generale con le proprietà indicate sotto:
\[T(n) \; = \; \begin{cases}
	\; T_0(n) & \text{se } n \leq n_0 \quad - \quad \text{caso base} \\[3pt]
	\; s(n) \cdot T(f(n)) + \omega(n) & \text{se } n > n_0 \quad - \quad \text{passo ricorsivo}
\end{cases}\]
\begin{itemize}
	\item \(T_0(n)\) è il costo dell'algoritmo diretto per risolvere un'istanza di taglia \(n \leq n_0\)
	\item \(s(n)\) è il numero di chiamate ricorsive effettuate per risolvere un'istanza di taglia \(n\) e coincide con
	il numero di nodi figli di ogni nodo dell'albero della ricorsione
	\item \(f(n)\) è la taglia della sottoistanza, generata dalla procedura di divide di un'istanza di taglia \(n\), che
	viene passata alla chiamata ricorsiva
	\item \(T(f(n))\) è il costo di ogni chiamata ricorsiva per risolvere un'istanza di taglia \(f(n)\)
	\item \(\omega(n) = T_D(n) + T_C(n)\) è il costo complessivo delle procedure di divide e di combine
\end{itemize}

\vspace{5pt}

\noindent
Per ottenere il valore analitico ed eliminare le ricorrenze è necessario risolvere l'equazione attraverso un'analisi
analitica, o ricorrendo al formulone o al master theorem. 

\subsubsection*{Iterate di \(f\)}
Si definisce l'iterata di una funzione \(f\) come la potenza della composizione di \(f\) con se stessa.
\[\begin{cases}
	f^{(0)}(n) = n \\[3pt]
	f^{(l)}(n) = f(f^{(l-1)})(n) & \text{per } l > 0
\end{cases}\]
Si osserva che \(f^{(i)}(n)\) corrisponde alla taglia dell'istanza all'i-esima chiamata ricorsiva. In particolare,
inizialmente \(f^{(i)}(n) = n\) per \(i = 0\) e progressivamente diminuisce ad ogni iterata, fino a raggiungere il caso
base \(f^{(k+1)}(n) \leq n_0\) per un determinato valore di \(i = k\). Si definisce \(k\) come il massimo valore di \(i\)
per cui la taglia dell'istanza è ancora maggiore di \(n_0\), ovvero:
\[k = f^*(n, n_0) := \max \left\{ k : f^{(k)}(n) > n_0 \right\}\]

\subsubsection*{Formulone}
Il formulone è una formula generale che risolve l'equazione generale alle ricorrenze indicata sopra, se sono soddisfatte
opportune condizioni indicate sotto.
\[
T(n) = \underbrace{
	\overbrace{\sum_{l=0}^{f^* (n, \, n_0)}}^{\begin{array}{c} \text{\footnotesize per ogni} \\ \text{\footnotesize livello interno} \end{array}}
	\overbrace{\left[ \prod_{j=0}^{l-1} s \left(f^{(j)}(n) \right) \right. \!}^{\begin{array}{c} \text{\footnotesize \# nodi al} \\ \text{\footnotesize livello \(l\)-esimo} \end{array}}
	\; \cdot \!\!\! \overbrace{\left. \phantom{\prod_{j}^{l} \!\!\!\!\!\!\!\!} \omega \left( f^{(l)}(n) \right) \right]}^{\begin{array}{c} \text{\footnotesize costo interno} \\ \text{\footnotesize di ogni nodo} \\ \text{\footnotesize al livello \(l\)-esimo} \end{array}}
	}_\text{lavoro compiuto nei nodi interni}
	\;\; + \;\;\;
	\underbrace{ \overbrace{\prod_{j=0}^{f^*(n, n_0)} s \left(f^{(j)}(n) \right) \!}^{\begin{array}{c} \text{\footnotesize \# di foglie} \end{array}}
	\; \cdot \; \overbrace{ \phantom{\prod_{j}^{f^*} \!\!\!\!\!\!\!\!} T_0 \left( f^{(f^*(n, n_0) + 1)}(n) \right)}^{\begin{array}{c} \text{\footnotesize costo di ogni foglia}\end{array}}
	}_\text{lavoro compiuto nelle foglie}
\]

\vspace{5pt}

\begin{itemize}
	\item il numero di figli dei nodi di un certo livello è costante, ovvero \(s(f^{(l)}(n)) = s_l\) per ogni livello \(l\)
	\item la taglia di tutti i nodi di un certo livello è costante, ovvero \(f^{(l)}(n) = f_l\) per ogni  livello \(l\)
\end{itemize}

\subsubsection*{Formulone applicato ad un caso particolare}
Di seguito un esempio di applicazione del formulone ad un caso particolare con:
\begin{itemize}
	\item \(s(n) = a\): numero di figli costante per ogni nodo interno
	\item \(f(n) = n/2\): la taglia dell'istanza si dimezza ad ogni chiamata ricorsiva
	\item \(\omega(n) = \beta \, n\): costo di divide and conquer lineare per ogni nodo interno
	\item \(T_0(n) = 1\): costo costante per ogni foglia
\end{itemize}
\begin{align*}
	T(n) \;\; &= \sum_{l=0}^{(\log_2 n) - 1} \left( a^l \cdot \beta \; \frac{n}{2^l} \right) \; + \; a^{\log_2 n} \cdot T_0(n) \; = \; \beta \, n \cdot \! \sum_{l=0}^{(\log_2 n) - 1} \!\! \left( a/2 \right)^l \; + \; n^{\log_2 a} = \\[3pt]
	&= \;\; \beta \, n \cdot \frac{(a/2)^{\log_2 n} - 1}{(a/2) - 1} + n^{\log_2 a} \; = \; \frac{\beta}{a/2 - 1} \cdot n \left( n^{\log_2 a/2} - 1\right) + n^{\log_2 a} = \\[3pt]
	&= \; \underbrace{\frac{\beta}{a/2 - 1} \cdot (n^{\log_2 a} - n)}_{\text{costo dei nodi interni}} \;\; + \!\! \underbrace{ \;\; n^{\log_2 a} \;\;}_{\text{costo delle foglie}}
	\!\! = \;\; \Theta \left( n^{\log_2 a} \right)
\end{align*}

\subsubsection*{Complessità degli algoritmi di ordinamento basati sui confronti}
Al momento non si conoscono funzioni che legano problemi alla relativa complessità asintotica e più in generale non si
conoscono metodi per individuare in maniera assoluta i lower bound di problemi, ad eccezione di quelli deducibili banalmente.

Ad esempio gli algoritmi di ordinamento basati sui confronti hanno lower bound banale di \(\Theta(n \log n)\), siccome
\(n \log n\) è la quantità di informazione necessaria a rappresentare tutte le possibili permutazioni di \(n\) elementi ed
equivale al numero minimo di confronti (rappresentabili in bit) necessari per distinguere tutte le possibili permutazioni.
\[\# \text{permutazioni} = n! \approx n^n \qquad\quad \text{informazione}(\text{perm} \,_i) = \log_2(n!) \approx \log_2(n^n) = n \log_2 n\]


\subsection{Master theorem}
\subsubsection*{Assunzioni iniziali}
Il master theorem è un teorema che fornisce una soluzione asintotica immediata in funzione di due parametri \(a\) e \(b\)
per una classe di equazioni alle ricorrenze con le seguenti caratteristiche:
\begin{itemize}
	\item \(s(n) = a\): numero di figli costante per ogni nodo interno
	\item \(f(n) = n/b\): la taglia dell'istanza si divide per una costante \(b > 1\) ad ogni chiamata ricorsiva
	\item \(n_0 = 1\): caso base per istanze di taglia 1
\end{itemize}

\[T(n) \; = \; \begin{cases}
	\; T_0(n) & \text{se } n \leq 1 \\[3pt]
	\; a \cdot T(n/b) + \omega(n) & \text{se } n > 1
\end{cases}\]

\subsubsection*{Funzione soglia}
Si definisce una \textbf{funzione soglia} \(\sigma(n)\) che rappresenta il numero di foglie dell'albero della ricorsione.
Permette di calcolare in maniera immediata il costo complessivo dovuto al lavoro compiuto nelle foglie che funge da lower
bound della complessità dell'algoritmo.
\[\sigma(n) = n^{\log_b a} \qquad T(n) \, \geq \, T_{\text{foglie}}(n) \, = \, T_0(n) \cdot \# \; \text{foglie} \, = \, T_0(n) \cdot \sigma(n)\]

\subsubsection*{Classi di funzioni e relative complessità asintotiche}
Paragonando il costo dovuto al lavoro compiuto da ogni nodo interno \(\omega(n)\) con la funzione soglia \(\sigma(n)\), si
individuano tre casi per cui è possibile calcolare la complessità asintotica di \(T(n)\) in maniera immediata:
\[T(n) \; = \; \begin{cases}
	\; \Theta (n^{\log_b a}) & \text{\footnotesize se il lavoro nei nodi interni \(\omega(n)\) è trascurabile rispetto a quello delle foglie \(\sigma(n)\)}  \\[3pt]
	\; \Theta (n^{\log_b a} \, \log_b n) & \text{\footnotesize se il lavoro nei nodi interni \(\omega(n)\) è dello stesso ordine di quello delle foglie \(\sigma(n)\)}  \\[3pt]
	\; \Theta (\omega(n)) & \text{\footnotesize se il lavoro nei nodi interni \(\omega(n)\) domina quello delle foglie \(\sigma(n)\)}
\end{cases}\]

\subsubsection*{Analisi del caso 1}
Si verifica quando il lavoro nei nodi interni \(\omega(n)\) è trascurabile rispetto a quello delle foglie \(\sigma(n)\):
\[\exists \varepsilon > 0 \text{ tale che } \omega(n) = O \! \left(n^{(\log_b a )- \epsilon}\right) = \sigma(n) \cdot O \! \left(n^{-\varepsilon}\right)\]
Di seguito si osserva che calcolando il contributo complessivo dei nodi interni usando \(\omega(n) = n ^{(\log_b a) - \varepsilon}\),
il risultato è comparabile con quello delle foglie senza far aumentare la complessità asintotica di \(\Theta(n^{\log_b a})\).
\begin{align*}
	T_{\text{int}}(n) &= \sum_{l=0}^{(\log_b n) - 1} a^l \cdot \omega \! \left( \frac{n}{b^l} \right) \; = \; \sum_{l=0}^{(\log_b n) - 1} a^l \cdot \left( \frac{n}{b^l} \right)^{(\log_b a) - \varepsilon} \; = \; n^{\log_b a} \; n^{-\varepsilon} \; \cdot \!\! \sum_{l=0}^{(\log_b n) - 1} \frac{a^l}{b^{l (\log_b a) - l \varepsilon}} \; = \\[3pt]
	&= \;\; n^{\log_b a} \; n^{-\varepsilon} \; \cdot \!\! \sum_{l=0}^{(\log_b n) - 1} \frac{a^l}{a^l \cdot b^{l \varepsilon}} \; = \; n^{\log_b a} \; n^{-\varepsilon} \; \cdot \!\! \sum_{l=0}^{(\log_b n) - 1} {(b^{\varepsilon})}^l \; = \; n^{\log_b a} \; n^{-\varepsilon} \cdot \frac{(b^{\varepsilon})^{(\log_b n)} - 1}{b^{\varepsilon} - 1} \; = \\[3pt]
	&= \;\; n^{\log_b a} \; n^{-\varepsilon} \cdot \frac{n^{\varepsilon} - 1}{b^{\varepsilon} - 1} \; \approx \; n^{\log_b a} \; n^{-\varepsilon} \cdot \frac{n^{\varepsilon}}{b^{\varepsilon} - 1} \; = \; \frac{n^{\log_b a}}{b^{\varepsilon} - 1} \; = \; \frac{\sigma(n)}{b^{\varepsilon} - 1} \; = \; O(\sigma(n))
\end{align*}
Si osserva che per \(\varepsilon \to 0^+\) il coefficiente moltiplicativo sarebbe molto elevato, ma comunque costante, per
cui asintoticamente trascurabile.

\subsubsection*{Analisi del caso 2}
Si verifica quando il lavoro nei nodi interni \(\omega(n)\) è dello stesso ordine di quello delle foglie \(\sigma(n)\):
\[\omega(n) = \Theta(n^{\log_b a}) = \Theta (\sigma(n))\]
Di seguito si dimostra che il contributo complessivo dei nodi interni ha una complessità asintotica maggiore di quella delle
foglie e domina la complessità totale dell'algoritmo, che risulta essere \(\Theta(n^{\log_b a} \, \log_b n)\).
\begin{align*}
	T_{\text{int}}(n) &= \sum_{l=0}^{(\log_b n) - 1} \left( a^l \cdot \omega \! \left( \frac{n}{b^l} \right) \right) \; = \; \sum_{l=0}^{(\log_b n) - 1} \left( a^l \cdot \left( \frac{n}{b^l} \right)^{\log_b a} \right) \; = \; n^{\log_b a} \; \cdot \!\! \sum_{l=0}^{(\log_b n) - 1} \left( \frac{a^l}{b^{l (\log_b a)}} \right) \\[3pt]
	&= n^{\log_b a} \; \cdot \!\! \sum_{l=0}^{(\log_b n) - 1} \frac{a^l}{a^l} \; = \; n^{\log_b a} \cdot (\log_b n) \; = \; \Theta(n^{\log_b a} \, \log_b n)
\end{align*}
Di seguito si osserva che il costo complessivo di ogni livello ha la stessa complessità asintotica di quello delle foglie,
per cui è possibile calcolare la complessità totale moltiplicando la complessità di ogni livello per il numero di livelli.
Gli algoritmi di questa tipologia (es. mergesort e quicksort) hanno la particolarità che il costo delle fasi di divide o
combine è proporzionale alla taglia dell'istanza.
\[T_{\text{\(l\)-esimo livello}}(n) = a^l \cdot \omega \! \left( \frac{n}{b^l} \right) = a^l \cdot \left( \frac{n}{b^l} \right)^{\log_b a} = a^l \cdot \frac{n^{\log_b a}}{b^{l (\log_b a)}} = a^l \cdot \frac{n^{\log_b a}}{a^l} = n^{\log_b a} = \sigma(n)\]

\newpage

\subsubsection*{Analisi del caso 3}
Si verifica quando il lavoro nei nodi interni \(\omega(n)\) domina quello delle foglie \(\sigma(n)\):
\[\exists \varepsilon > 0 \text{ tale che } \omega(n) = \Omega(n^{\log_b a + \epsilon}) = \sigma(n) \cdot \Omega(n^{\epsilon})\]
Dalla condizione precedente si può dedurre anche la seguente condizione di regolarità, ovvero che il lavoro interno complessivo
nei nodi figli (al livello inferiore) è minore del lavoro interno del nodo padre.
\[\exists c \in \mathbb{R} \text{ con } 0 < c < 1 \text{ tale che } \forall n \in \mathbb{N} \text{ si ha } a \cdot \omega \! \left( \frac{n}{b} \right) \leq c \cdot \omega(n)\]
Per la dimostrazione del terzo caso si utilizza la condizione di regolarità appena enunciata, in quanto rende lo svolgimento
più semplice e immediato. Di seguito sono riportati alcuni passaggi preliminari che verranno poi usati nella successiva
dimostrazione.
\[a\cdot \omega \! \left(\frac{n}{b}\right) \leq c\cdot \omega(n) \;\; \Rightarrow \;\; \omega \! \left(\frac{n}{b}\right) \leq \frac{c}{a} \cdot \omega(n) \;\; \Rightarrow \;\; \omega \! \left(\frac{n}{b^l}\right) \leq \left( \frac{c}{a} \right)^l \cdot \omega(n)\]
\[T_{\text{int}}(n) = \!\! \sum_{l=0}^{(\log_b n) - 1} a^l \cdot \omega \! \left( \frac{n}{b^l} \right) \; \leq \!\! \sum_{l=0}^{(\log_b n) - 1} a^l \cdot \frac{c^l}{a^l} \cdot \omega(n) \; = \; \omega(n) \cdot \!\!\! \sum_{l=0}^{(\log_b n) - 1} c^l \; < \; \omega(n) \cdot \frac{1}{1 - c} \; = \; \Theta(\omega(n))\]
Dall'analisi precedente si osserva che il costo complessivo dei nodi interni è proporzionale al costo di ogni nodo interno
e domina la complessità totale dell'algoritmo se il costo del singolo nodo interno è asintoticamente maggiore di quello delle
foglie, ovvero se \(\omega(n) = \Omega(n^{\log_b a + \epsilon})\).

\subsubsection*{Considerazione sugli altri casi}
Esistono funzioni che non possono essere classificate in nessuno dei tre casi del master theorem, ad esempio quando
\(\omega(n) = n^{\log_b a} \log n\). In questi casi è necessario risolvere l'equazione alle ricorrenze attraverso
altri metodi risolutivi, ad esempio il formulone o l'analisi analitica.

\subsection{Applicazioni del divide and conquer}
Nei seguenti capitoli sono riportati alcuni esempi di applicazioni efficienti del paradigma divide and conquer:
\begin{itemize}
	\item esponenziazione con esponente intero
	\item moltiplicazione di polinomi tramite l'algoritmo di Karatsuba
	\item moltiplicazione di matrici tramite l'algoritmo di Strassen
	\item moltiplicazione di matrici supercircolanti tramite la trasformata di Hadamard
	\item moltiplicazione di polinomi tramite la trasformata di Fourier (FFT) e l'algoritmo di Cooley-Tukey
\end{itemize}

\newpage

\section{Esponenziazione con esponente intero}
\subsubsection*{Definizione formale del problema}
Dato un elemento \(x\) proveniente da una struttura algebrica con operazione associativa (ad esempio numeri in \(\mathbb{N}\),
\(\mathbb{R}\), \(\mathbb{C}\), matrici, \dots) costante, si vuole calcolare \(x^n\) con \(n \in \mathbb{N}\). Siccome si
tratta \(x\) come costante, si esclude che possa essere un polinomio per semplificare la risoluzione.

\subsubsection*{Osservazioni sulla taglia del problema}
Essendo \(x\) costante, viene escluso dal calcolo della taglia che è determinata unicamente da \(n\). Si osserva che la
quantità di informazione necessaria a rappresentare \(n\) è \(\log_2 n\), per cui la taglia del problema in funzione di
cui va espressa la complessità è \(\text{taglia}(n) = |n| = \log_2 n\), con \(n = 2 ^{\text{\, taglia}(n)}\).

\subsection{Algoritmo ovvio basato sulla definizione}
L'algoritmo ovvio per calcolare \(x^n\) si basa sulla definizione ricorsiva di potenza come composizione multipla
dell'operazione associativa della struttura algebrica, ovvero:
\[x^n := \begin{cases}
	\; x^{(0)} = 1 & \text{se } n = 0 \\[3pt]
	\; x^{(n-1)} \; \circ \; x & \text{se } n > 0 \quad - \quad \circ \;\; \text{è l'operazione associativa della struttura algebrica}
\end{cases}\]
La complessità equivale al costo di \(n - 1\) composizioni dell'operazione associativa, per cui la complessità risulta
essere esponenziale in funzione della taglia del problema, ovvero:
\[T(n) = n - 1 = 2^{\, \text{taglia}(n)} - 1 = \Theta \left( 2^{\, \text{taglia}(n)} \right)\]

\subsection{Algoritmo efficiente basato sul divide and conquer}
Si osserva che se l'operazione è associativa, come richiesto nelle ipotesi, è possibile riarrangiare le composizioni
di funzioni in modo da riutilizzare i risultati di composizioni già calcolate. Di seguito la definizione ricorsiva
dell'algoritmo efficiente basato sul divide and conquer:
\[x^n := \begin{cases}
	\; x^{(0)} = 1 & \text{se } n = 0 \\[3pt]
	\; x^{(1)} = x & \text{se } n = 1 \\[3pt]
	\; x^{(n/2)} \; \circ \; x^{(n/2)} & \text{se } n > 1 \text{ è pari} \\[3pt]
	\; x^{\lfloor n/2 \rfloor} \; \circ \; x^{\lfloor n/2 \rfloor} \; \circ \; x & \text{se } n > 1 \text{ è dispari}
\end{cases}\]
L'equazione alle ricorrenze per questo algoritmo al caso peggiore è la seguente, osservando che è sufficiente una chiamata
ricorsiva per calcolare \(x^{(n/2)}\) e, al caso peggiore, si aggiungono 2 operazioni di composizione.
\[T(n) = \begin{cases}
	\; 1 & \text{se } n = 0, 1 \\[3pt]
	\; T(n/2) + 2 & \text{se } n > 1
\end{cases}\]
Utilizzando il master theorem con \(a = 1\), \(b = 2\) e \(\omega(n) = O(1)\), si ottiene una complessità lineare
in funzione della taglia del problema, ovvero:
\[T(n) = \Theta \left( \log_2 n \right) = \Theta \left( \text{taglia}(n) \right)\]

\newpage

\section{Moltiplicazione di polinomi - Karatsuba}
\subsubsection*{Definizione formale di polinomio e degree bound}
Un polinomio \(P(x)\) di grado \(N-1\) (con \(N\) termini) è definito come:
\[P(x) \; := \; \sum_{i=0}^{N-1} \, a_i \, x^i \qquad\qquad \text{con} \;\; a_0, a_1, \dotsc \; a_{N-1} \; \text{fissati}\]
Si osserva che un polinomio può essere visto anche come un numero espresso in base \(x\) con \(0\leq a_i < x\) secondo la
notazione posizionale.

Il degree bound \(\alpha\) è un limite superiore al grado di un polinomio, che indica l'insieme di tutti i polinomi di
grado inferiore o uguale ad \(\alpha\), ovvero \(N - 1 \leq \alpha\).

\subsection{Prodotto tra polinomi usando la definizione}
\subsubsection*{Tecnica risolutiva}
Dati due polinomi \(P(x)\) e \(Q(x)\) di grado \(N - 1\), il loro prodotto \(P(x) \cdot Q(x)\) è definito come segue.
\[P(x) \; := \; \sum_{i=0}^{N-1} \, a_i \, x^i \qquad\quad Q(x) \; := \; \sum_{i=0}^{N-1} \, b_i \, x^i \qquad\quad P(x) \cdot Q(x) \; = \; \sum_{i=0}^{2N-2} \, c_i \, x^i\]

\subsubsection*{Analisi della complessità}
Per calcolare il prodotto usando la definizione sono necessarie:
\begin{itemize}
	\item \(N^2\) moltiplicazioni tra i coefficienti \(a_i\) e \(b_j\)
	\item \((N-1)^2 = N^2 - 1 - (2N - 2)\) addizioni tra i prodotti \(a_i \cdot b_j\) per ottenere i coefficienti \(c_k\)
\end{itemize}
Complessivamente il prodotto usando la definizione ha complessità \(\Theta(N^2)\).

\subsection{Prodotto tra polinomi con approccio divide and conquer}
\subsubsection*{Tecnica risolutiva}
I polinomi \(P(x)\) e \(Q(x)\), di grado \(N-1\), possono essere scomposti in polinomi più semplici di grado \(< N/2\)
e il prodotto tra \(P(x)\) e \(Q(x)\) risulta essere in funzione del prodotto di questi polinomi più semplici, ovvero:
\[P(x) = A(x) + B(x) \cdot x^{N/2} \qquad\qquad Q(x) = C(x) + D(x) \cdot x^{N/2}\]
\[P(x) \cdot Q(x) = B(x) \, D(x) \cdot x^N + (A(x) \, D(x) + B(x) \, C(x)) \cdot x^{N/2} + A(x) \, C(x)\]

\subsubsection*{Equazione alle ricorrenze}
L'equazione alle ricorrenze per questo approccio è la seguente:
\[T(N) \; = \; \begin{cases}
	\; 1 & \text{se } N = 1 \\[3pt]
	\; 4 \cdot T(N/2) + \beta_{std} \; N & \text{se } N > 1
\end{cases}\]
\begin{itemize}
	\item il caso base è rappresentato dal prodotto di due polinomi di grado 0 che richiede un solo prodotto
	\item il passo ricorsivo prevede 4 chiamate ricorsive per calcolare i quattro prodotti \(AC\), \(AD\), \(BC\) e \(BD\)
	tra polinomi di grado \(< N/2\)
	\item il costo per ogni nodo interno è dato dalle somme dei coefficienti dei termini di grado uguale, che è proporzionale
	alla taglia \(N\)
\end{itemize}

\subsubsection*{Analisi della complessità}
Dal formulone si ottiene che la complessità rimane sempre \(\Theta(N^2)\), da come evidenziato di seguito:
\begin{align*}
	T(N) \;\; &= \sum_{l=0}^{(\log_2 N) - 1} \!\! \left( 4^l \cdot \beta_{std} \; \frac{N}{2^l} \right) + 4^{\log_2 N} \cdot T_0(n) \; = \; \beta_{std} \; N \cdot \! \sum_{l=0}^{(\log_2 N) - 1} \!\! \left( 2^{2l}/2^l \right) + N^{\log_2 4} \; = \\[3pt]
	&= \;\; \beta_{std} \; N \cdot \sum_{l=0}^{(\log_2 N) - 1} \!\! 2^l + N^2 \; = \; \beta_{std} \; N \cdot \frac{2^{\log_2 N} - 1}{2 - 1} + N^2 \; = \\[3pt]
	&= \, \underbrace{\, \beta_{std} \; (N^2 - N) \,}_{\begin{array}{c} \text{\scriptsize costo delle somme} \\[-3pt] \text{ \scriptsize nei nodi interni} \end{array}}
	\;\; + \!\!\!\!\!\! \underbrace{\;\;\;N^2\;\;}_{\begin{array}{c} \text{\scriptsize costo dei prodotti} \\[-3pt] \text{\scriptsize nelle foglie} \end{array}} \!\!\!\!\!\! = \;\; \Theta \, (N^2)
\end{align*}

\subsection{Prodotto di polinomi con l'algoritmo di Karatsuba (1962)}
\subsubsection*{Tecnica risolutiva}
Karatsuba osserva che è possibile ricavare la somma dei due prodotti \(AD + BC\) riutilizzando i due prodotti \(AC\) e \(BD\)
già precedentemente calcolati come segue:
\[(A + B) \cdot (C + D) = AC + (AD + BC) + BD \implies AD + BC = (A + B) \cdot (C + D) - AC - BD\]
In questo modo il numero di prodotti da calcolare ricorsivamente si riduce da 4 a 3, con l'effetto di aumentare lievemente
il costo della fase di divide dovuto al calcolo delle somme \((A + B)\) e \((C + D)\).

\subsubsection*{Equazione alle ricorrenze ed analisi della complessità}
L'equazione alle ricorrenze per l'algoritmo di Karatsuba è la seguente:
\[T(N) \; = \; \begin{cases}
	\; 1 & \text{se } N = 1 \\[3pt]
	\; 3 \cdot T(N/2) + \beta_{kz} \; N & \text{se } N > 1
\end{cases}\]

Dal formulone si ottiene che la complessità è \(\Theta(N^{\log_2 3}) \approx \Theta(N^{1.58})\), da come evidenziato di seguito:
\begin{align*}
	T(N) \;\; &= \sum_{l=0}^{(\log_2 N) - 1} \!\! \left( 3^l \cdot \beta_{kz} \; N / 2^l \right) + 3^{\log_2 N} \cdot T(1) \; =  \; \beta_{kz} \; N \cdot \! \sum_{l=0}^{(\log_2 N) - 1} \!\! \left( (3/2)^l \right) + N^{\log_2 3} \; = \\[3pt]
	&= \;\; \beta_{kz} \; N \cdot \frac{(3/2)^{\log_2 N} - 1}{(3/2) - 1} + N^{\log_2 3} \; = \; 2 \beta_{kz} \; N \cdot (N^{\log_2 3 - \log_2 2} - 1) + N^{\log_2 3} \; = \\[3pt]
	&= \;\; \underbrace{\, 2 \beta_{kz} \; (N^{\log_2 3} - N) \,}_{\begin{array}{c} \text{\scriptsize costo delle somme} \\[-3pt] \text{\scriptsize nei nodi interni} \end{array}}
	\;\; + \!\!\!\!\!\! \underbrace{\;\;N^{\log_2 3}\;\;}_{\begin{array}{c} \text{\scriptsize costo dei prodotti} \\[-3pt] \text{\scriptsize nelle foglie} \end{array}} \!\!\!\!\!\! = \;\; \Theta \, (N^{\log_2 3}) \; \approx \; \Theta \, (N^{1.58})
\end{align*}

\subsubsection*{Considerazioni sull'algoritmo di Karatsuba}
\begin{itemize}
	\item per valori di \(N\) piccoli, l'algoritmo standard (D\&C) rimane ancora più efficiente di quello di Karatsuba,
	in quanto il costo aggiuntivo della fase di divide è maggiore del risparmio ottenuto dalla chiamata ricorsiva in meno
	\item per valori di \(N\) molto grandi, è possibile migliorare ulteriormente l'efficienza asintotica nel calcolo del
	prodotto di polinomi utilizzando ad esempio l'algoritmo di Toom-Cook
	\item la vera novità dell'algoritmo di Karatsuba è stata quella di trovare il primo algoritmo sub-quadratico per
	risolvere un problema che tutti i più grandi matematici ritenevano impossibile da risolvere in meno di \(\Theta(N^2)\)
\end{itemize}

\newpage

\section{Moltiplicazione di matrici - Strassen}
\subsubsection*{Note introduttive sulle matrici}
Per la seguente sezione relativa al prodotto di matrici, si considerano solo matrici quadrate \(n \times n\)
(\(\in \mathbb{R}^{n \times n}\)) con \(n = 2^k\) potenza di 2, per un certo \(k \in \mathbb{N}\). In questo modo
le matrici sono perfettamente scomponibili ricorsivamente in quattro sottomatrici di dimensione \(n/2 \times n/2\).

Per convenzione si assume che la taglia di una matrice è pari al numero di righe (o di colonne) di cui è composta,
ovvero \(A \in \mathbb{R}^{n \times n} \;\Rightarrow\; |A| = n\).

\subsection{Prodotto di matrici usando la definizione}
\subsubsection*{Tecnica risolutiva}
Date due matrici \(A, B \in \mathbb{R}^{n \times n}\), con \(n = 2^k\), il loro prodotto \(C = A \cdot B\) è sempre una
matrice \(C \in \mathbb{R}^{n \times n}\) i cui elementi \(c_{ij}\) sono calcolati come segue:
\[c_{ij} = \sum_{k=1}^{n} a_{ik} \cdot b_{kj} \qquad\qquad \text{per } i, j = 1, 2, \dotsc, n\]
In altre parole il prodotto di matrici è una matrice che ha come elementi il prodotto interno dei vettori riga e colonna
rispettivamente della prima matrice \(A\) e della seconda matrice \(B\).

\subsubsection*{Analisi della complessità}
Il prodotto interno tra un vettore riga e un vettore colonna richiede \(n\) moltiplicazioni e \(n-1\) addizioni, per un costo
totale di \(T(n) = 2n - 1\). Siccome vanno fatti \(n^2\) prodotti interni, il costo totale della moltiplicazione di matrici
risulta \(\Theta(n^3)\): \[T(n) = n^2 (2n - 1) = 2n^3 - n^2 = \Theta(n^3)\]

\subsection{Prodotto di matrici con approccio divide and conquer}
\subsubsection*{Tecnica risolutiva}
Le matrici \(A, B \in \mathbb{R}^{n \times n}\) possono essere scomposte ciascuna in quattro sottomatrici di dimensioni
\(n/2 \times n/2\), identificati con \(X_{11}, X_{12}, X_{21}, X_{22} \in \mathbb{R}^{n/2 \times n/2}\). Il prodotto
\(C = A \cdot B\) risulta essere in funzione dei prodotti tra queste sottomatrici.
\[\begin{array}{c c}
	C_{11} = A_{11} \cdot B_{11} + A_{12} \cdot B_{21} \qquad\quad & \qquad\quad C_{21} = A_{21} \cdot B_{11} + A_{22} \cdot B_{21} \\[5pt]
	C_{12} = A_{11} \cdot B_{12} + A_{12} \cdot B_{22} \qquad\quad & \qquad\quad C_{22} = A_{21} \cdot B_{12} + A_{22} \cdot B_{22}
\end{array}\]

\subsubsection*{Equazione alle ricorrenze}
L'equazione alle ricorrenze per questo approccio è la seguente:
\[T(n) \; = \; \begin{cases}
	\; 1 & \text{se } n = 1 \\[3pt]
	\; 8 \cdot T(n/2) + 4 n^2 & \text{se } n > 1
\end{cases}\]
\begin{itemize}
	\item il caso base è rappresentato dal prodotto di due matrici \(1 \times 1\) che richiede un solo prodotto
	\item il passo ricorsivo prevede 8 chiamate ricorsive per calcolare gli otto prodotti tra le sottomatrici di dimensione \(n/2 \times n/2\)
	\item il costo per ogni nodo interno è dato dalle 4 somme dei prodotti tra sottomatrici \(n/2 \times n/2\)
\end{itemize}

\subsubsection*{Analisi della complessità}
Dal formulone si ottiene che la complessità rimane sempre \(\Theta(n^3)\), da come evidenziato di seguito:
\begin{align*}
	T(n) \;\; &= \sum_{l=0}^{(\log_2 n) - 1} \!\! \left( 8^l \cdot \left( \frac{n}{2^l} \right)^2 \right) + 8^{\log_2 n} \cdot T_0(n) \; = \; n^2 \cdot \! \sum_{l=0}^{(\log_2 n) - 1} \!\! 2^l + n^{\log_2 8} = n^2 \cdot \frac{2^{\log_2 n} - 1}{2 - 1} + n^3 \; = \\[3pt]
	&= \underbrace{\;n^2 \cdot (n - 1)\;}_{\begin{array}{c} \text{\scriptsize costo delle somme} \\[-3pt] \text{\scriptsize dei nodi interni} \end{array}}
	+ \!\!\!\!\!\! \underbrace{\;\;\; n^3 \;\;}_{\begin{array}{c} \text{\scriptsize costo dei prodotti} \\[-3pt] \text{\scriptsize alle foglie} \end{array}}
	\!\!\!\!\!\! =  \;\; 2n^3 - n^2  \; = \; \Theta(n^3)
\end{align*}

\subsubsection*{Vantaggi dell'algoritmo D\&C ricorsivo rispetto a quello iterativo basato sulla definizione}
Ad oggi le macchine moderne performano molto meglio con l'algoritmo D\&C ricorsivo rispetto a quello iterativo basato sulla
definizione in quanto si lavora con sottomatrici sempre più piccole fino a quando riescono ad essere totalmente memorizzate
all'interno della cache. In questo modo si riducono drasticamente i tempi di accesso alla memoria principale dovuti ai
continui cache miss quando si scorrono per intero righe e colonne per calcolare i prodotti interni.

\subsection{Prodotto di matrici con l'algoritmo di Strassen (1969)}
\subsubsection*{Tecnica risolutiva}
Ispirandosi a Karatsuba, Strassen osserva che è possibile ricavare gli 8 prodotti tra le sottomatrici \(A_{ij}\) e \(B_{ij}\)
riutilizzando solo 7 prodotti, calcolati su opportune somme e differenze di sottomatrici.
\vspace{-12pt}
\begin{center}
	\begin{minipage}{0.25\textwidth}
		\begin{align*}
			A_1 &= (A_{12} - A_{22}) \\
			A_2 &= (A_{11} + A_{22}) \\
			A_3 &= (A_{11} - A_{21}) \\
			A_4 &= (A_{11} + A_{12}) \\
			A_5 &= A_{11} \\
			A_6 &= A_{22} \\
			A_7 &= (A_{21} + A_{22})
		\end{align*}
	\end{minipage}
	\begin{minipage}{0.25\textwidth}
		\begin{align*}
			B_1 &= (B_{21} + B_{22}) \\
			B_2 &= (B_{11} + B_{22}) \\
			B_3 &= (B_{11} + B_{12}) \\
			B_4 &= B_{22} \\
			B_5 &= (B_{12} - B_{22}) \\
			B_6 &= (B_{21} - B_{11}) \\
			B_7 &= B_{11}
		\end{align*}
	\end{minipage}
	\begin{minipage}{0.4\textwidth}
		\begin{align*}
			P_i = A_i \cdot B_i \qquad \text{per } i = 1, 2, \dotsc, 7
		\end{align*}
		\vspace{-20pt}
		\begin{align*}
			C_{11} &= P_1 + P_2 - P_4 + P_6 \\
			C_{12} &= P_4 + P_5 \\
			C_{21} &= P_6 + P_7 \\
			C_{22} &= P_2 - P_3 + P_5 - P_7
		\end{align*}
	\end{minipage}
\end{center}

\subsubsection*{Osservazioni sui prodotti di Strassen}
\begin{itemize}
	\item se si sviluppano i conti per le quattro espressioni di \(C_{ij}\) si può notare che si ottengono esattamente
	le stesse espressioni dell'approccio D\&C standard
	\item esiste una simmetria tra le espressioni dei blocchi \(A_i\) e \(B_i\): tra \(A_1\) e \(B_5\), tra \(A_2\)
	e \(B_2\), tra \(-A_3\) e \(B_6\), tra \(A_4\) e \(B_3\), tra \(A_5\) e \(B_7\), tra \(A_6\) e \(B_4\), tra \(A_7\) e \(B_1\)
\end{itemize}

\subsubsection*{Equazione alle ricorrenze}
L'equazione alle ricorrenze per l'algoritmo di Strassen è la seguente:
\[T(n) \; = \; \begin{cases}
	\; 1 & \text{se } n = 1 \\[3pt]
	\; 7 \cdot T(n/2) + 18 \; n^2/4 & \text{se } n > 1
\end{cases}\]

\begin{itemize}
	\item il caso base è rappresentato dal prodotto di due matrici \(1 \times 1\) che richiede un solo prodotto
	\item il passo ricorsivo prevede 7 chiamate ricorsive per calcolare i sette prodotti \(P_i\) tra i blocchi \(A_i\) e \(B_i\)
	di dimensione \(n/2 \times n/2\)
	\item il costo della fase di divide consiste nelle 10 somme e differenze tra sottomatrici \(A_{ij}, B_{ij} \in \mathbb{R}^{n/2 \times n/2}\)
	per calcolare i blocchi \(A_i\) e \(B_i\) (complessivamente \(10 \cdot n^2/4\))
	\item il costo della fase di conquer consiste nelle 8 somme/differenze tra i prodotti \(P_i \in \mathbb{R}^{n/2 \times n/2}\)
	per calcolare i blocchi \(C_{ij}\) (complessivamente \(8 \cdot n^2/4\))
\end{itemize}

\subsubsection*{Analisi della complessità}
Applicando il formulone si ottiene che la complessità è \(\Theta(n^{\log_2 7}) \approx \Theta(n^{2.81})\), da come mostrato di seguito:
\begin{align*}
	T(n) \;\; &= \sum_{l=0}^{(\log_2 n) - 1} \!\! \left( 7^l \cdot \frac{18}{4} \cdot \left( \frac{n}{2^l} \right)^2 \right) + 7^{\log_2 n} \cdot T_0(n) \; = \; \frac{9}{2} \; n^2 \cdot \! \sum_{l=0}^{(\log_2 n) - 1} \!\! \left( (7/4)^l \right) + n^{\log_2 7} \; = \\[3pt]
	&= \;\; \frac{9}{2} \; n^2 \cdot \frac{(7/4)^{\log_2 n} - 1}{(7/4) - 1} + n^{\log_2 7} \; = \; \frac{9}{2} n^2 \cdot \frac{n^{\log_2 (7/4)} - 1}{3/4} + n^{\log_2 7} \; = \\[3pt]
	&= \;\; \frac{9}{2} \cdot \frac{4}{3} \; n^2 \cdot \left( n^{(\log_2 7) - 2} - 1 \right) + n^{\log_2 7} \; = \; 6 n^2 \cdot (n^{\log_2 7}/n^2 - 1) + n^{\log_2 7} \\[3pt]
	&= \; \underbrace{\; 6 n^{\log_2 7} - 6n^2 \;}_{\begin{array}{c} \text{\scriptsize costo delle somme} \\[-3pt] \text{\scriptsize dei nodi interni} \end{array}}
	\;\; + \!\!\!\! \underbrace{\;\; n^{\log_2 7} \;\;}_{\begin{array}{c} \text{\scriptsize costo dei prodotti} \\[-3pt] \text{\scriptsize alle foglie} \end{array}}
	\!\!\!\!\!\! = \;\; 7n^{\log_2 7} - 6n^2 \; = \; \Theta(n^{\log_2 7}) \; \approx \; \Theta(n^{2.81})
\end{align*}

\subsubsection*{Osservazioni generali sull'algoritmo di Strassen}
\begin{itemize}
	\item una volta le moltiplicazioni di numeri erano molto costose, ora i blocchi logici sono altamente ottimizzati; inoltre
	molti processori dispongono dell'istruzione MADD (multiply and add) che consente di eseguire in un solo ciclo di clock
	una moltiplicazione seguita da un'addizione
	\item se si volesse eseguire il quadrato di una matrice \(A\) con l'algoritmo di Strassen, si risparmierebbero 5 somme
	e sottrazioni siccome \(A_i = B_i \;\; \forall i\)
	\item in generale non si può avere corrispondenza tra i blocchi \(A_i\) e \(B_i\), perché se così fosse, l'algoritmo
	diventerebbe commutativo portando sicuramente ad un risultato sbagliato siccome la moltiplicazione di matrici non è
	commutativa
	\item utilizzando la definizione, la complessità asintotica del prodotto di matrici era \(n^\Omega\) con \(\Omega = 3\),
	con Strassen si è riusciti a ridurre il valore di \(\Omega\) a \(\log_2 7 \approx 2.81\); ad oggi si conoscono algoritmi
	per il prodotto di matrici con \(\Omega \approx 2.37\), ma risultano poco vantaggiosi in quanto hanno costanti asintotiche
	molto elevate che li rendono svantaggiosi per valori di \(n\) piccoli
\end{itemize}

\subsection{Algoritmo ibrido per il prodotto di matrici}
\subsubsection*{Idea implementativa}
Per valori di \(n \leq n_0\) si osserva che l'algoritmo standard (D\&C) rimane ancora più efficiente di quello di Strassen.
Si vuole determinare tale valore \(n_0\) e sviluppare un algoritmo tale che per \(n > n_0\) utilizzi l'algoritmo di Strassen, 
mentre per \(n < n_0\) utilizzi l'algoritmo standard (D\&C).

\subsubsection*{Calcolo del parametro \(n_0\)}
Per calcolare il parametro \(n_0\), si cerca il valore che minimizza la complessità dell'algoritmo di Strassen imponendo
che per \(n \leq n_0\) venga utilizzato l'algoritmo standard (D\&C), ovvero:
\begin{itemize}
	\item i livelli dell'albero della ricorsione arrivano a \((\log_2 (n/n_0)) - 1\)
	\item il costo delle foglie è pari alla complessità dell'algoritmo standard (D\&C), ovvero \(T_0(n) = 2n^3 - n^2\)
\end{itemize}
\begin{align*}
	T(n, n_0) \;\; &= \sum_{l=0}^{(\log_2 (n/n_0)) - 1} \!\! \left( 7^l \cdot \frac{18}{4} \cdot \left( \frac{n}{2^l} \right)^2 \right) \; + \; 7^{\log_2 (n/n_0)} \cdot \left( 2n_0^3 - n_0^2 \right) \; = \; \dots \; = \\[3pt]
	&= \;\; n^{\log_2 7} \cdot \underbrace{\frac{2 n_0 + 5}{{n_0}^{(\log_2 7) - 2}}}_{g(n_0)} - 6n^2
\end{align*}
Si procede a minimizzare la funzione \(T(n, n_0)\) per trovare il valore ottimale di \(n_0\):
\vspace{-5pt}
\[\min_{n_0} \, T(n, n_0) \;\; \longrightarrow \;\; \min_{n_0} \, g(n_0) \;\; \longrightarrow \;\; \begin{array}{l} n_{0, \min} \approx 10.477 \\[3pt] g(n_{0, \min}) = 3.895 \end{array}\]

\subsubsection*{Considerazioni sul parametro \(n_0\)}
Dal punto di vista teorico, si sceglie \(n_0 = 8\) in quanto è la potenza di 2 più vicina a \(n_{0, \min}\). Riscrivendo
l'equazione alle ricorrenze e la complessità asintotica dell'algoritmo ibrido, si ottiene:
\[T(n, n_{0, \min}) \; = \; \begin{cases}
	\; 2n^3 - n^2 & \text{se } n \leq 8 \\[3pt]
	\; 7 \cdot T(n/2) + 18 \; n^2/4 & \text{se } n > 8
\end{cases} \qquad T(n, n_{0, \min}) \; = \begin{cases}
	\; 2n^3 - n^2 & \text{se } n \leq 8 \\[3pt]
	\; 3.918 \cdot n^{\log_2 7} - 6n^2 & \text{se } n > 8
\end{cases}\]

Dal punto di vista pratico, non è detto che \(n_0 = 8\) sia il valore ottimale, in quanto il modello teorico non tiene conto
di overhead dovuti al cambio di algoritmo, alla gestione della cache e altri fattori specifici della macchina su cui viene
eseguito l'algoritmo.

Ad oggi si sta sviluppando una branca di ricerca chiamata \textbf{adaptive software} che punta ad aggiustare dinamicamente i parametri
degli algoritmi di algebra lineare (come \(n_0\)) attraverso analisi e test pratici effettuati sulla macchina stessa, al fine
di ottenere le migliori prestazioni possibili per quella determinata macchina.

\newpage

\section{Moltiplicazione di matrici supercircolanti - Tr. di Hadamard}
\subsubsection*{Definizione di matrici supercircolanti}
Una matrice \(A \in \mathbb{C}^{n \times n}\), con \(n=2^k\) si dice supercircolante se:
\begin{itemize}
	\item per \(n = 1\) (banalmente tutti i numeri sono considerati matrici supercircolanti))
	\item per \(n > 1\) la matrice si può dividere in quattro blocchi \(n/2 \times n/2\) uguali a due a due lungo la diagonale
	principale e la diagonale secondaria e tali blocchi sono a loro volta supercircolanti
	\[A = \begin{bmatrix} A_1 & A_2 \\ A_2 & A_1 \end{bmatrix} \]
\end{itemize}

\subsubsection*{Osservazioni sulle matrici supercircolanti}
\begin{itemize}
	\item nelle matrici supercircolanti, ogni riga è ottenuta tramite uno shift circolare della riga precedente
	\item le righe successive alla prima sono univocamente definite a partire dalla prima riga, detta riga generatrice,
	tramite opportuni shift circolari
	\item dato che è sufficiente la prima riga per definire una matrice supercircolante, la taglia di una matrice
	supercircolante \(A\) è pari al numero di elementi della sua prima riga, ovvero \(n\)
	\item si osserva che l'insieme delle matrici supercircolanti è chiuso rispetto alla somma:
	\[A + B = \begin{bmatrix} A_1 & A_2 \\ A_2 & A_1 \end{bmatrix} + \begin{bmatrix} B_1 & B_2 \\ B_2 & B_1 \end{bmatrix} = \begin{bmatrix} A_1 + B_1 & A_2 + B_2 \\ A_2 + B_2 & A_1 + B_1 \end{bmatrix} \quad \text{supercircolante per def.}\]
	\item si osserva che l'insieme delle matrici supercircolanti è chiuso rispetto alla moltiplicazione:
	\[A \cdot B = \begin{bmatrix} A_1 & A_2 \\ A_2 & A_1 \end{bmatrix} \cdot \begin{bmatrix} B_1 & B_2 \\ B_2 & B_1 \end{bmatrix} = \begin{bmatrix} A_1 B_1 + A_2 B_2 & A_1 B_2 + A_2 B_1 \\ A_2 B_1 + A_1 B_2 & A_2 B_2 + A_1 B_1 \end{bmatrix} \quad \text{supercircolante per def.}\]
	\item siccome l'insieme delle matrici supercircolanti è chiuso rispetto alle operazioni di somma e moltiplicazione, allora
	è considerato un anello, in più si osserva che è anche commutativo, siccome le due operazioni godono della proprietà
	commutativa, ovvero \(A + B = B + A\) e \(A \cdot B = B \cdot A\)
\end{itemize}

\subsubsection*{Somme tra matrici supercircolanti}
Per sommare due matrici supercircolanti è sufficiente sommare le prime righe dei due addendi elemento per elemento, ottenendo
la prima riga della matrice risultante dalla somma. La complessità di questo approccio è \(T_{add}(n) = \Theta(n)\) con \(n\)
pari alla taglia delle matrici.

\subsection{Moltiplicazioni tra matrici supercircolanti}
\subsubsection*{Moltiplicazioni usando la definizione}
Usando la definizione di moltiplicazione di matrici, si osserva che si richiede il calcolo di 4 prodotti tra sottomatrici
supercircolanti di dimensione \(n/2 \times n/2\) e 2 somme tra matrici supercircolanti di dimensione \(n/2 \times n/2\).
L'equazione alle ricorrenze per questo approccio è la seguente:
\[T(n) \; = \; \begin{cases}
	\; 1 & \text{se } n = 1 \\[3pt]
	\; 4 \cdot T(n/2) + 2 \cdot n/2 & \text{se } n > 1
\end{cases} \;\; = \;\; \begin{cases}
	\; 1 & \text{se } n = 1 \\[3pt]
	\; 4 \cdot T(n/2) + n & \text{se } n > 1
\end{cases}\]
Dal master theorem con \(a = 4\), \(b = 2\) e \(\omega(n) = n\) si ottiene una funzione soglia di
\(\sigma(n) = n^{\log_2 4} = n^2\), tale per cui \(\omega(n) = \sigma(n) \cdot O(n^{-1})\), ovvero si ricade
nel primo caso del master theorem. La complessità è quindi \(T(n) = \Theta(\sigma(n)) = \Theta(n^2)\).

\newpage

\subsubsection*{Moltiplicazioni con approccio D\&C ottimizzato}
Si osserva che secondo opportune manipolazioni algebriche è possibile ricavare i due blocchi \(C_1\) e \(C_2\) del prodotto
attraverso soli 2 prodotti, con l'overhead di 6 somme e due moltiplicazioni per uno scalare.
\[\begin{array}{l}
	U = (A_1 + A_2) \cdot (B_1 + B_2) = A_1 B_1 + A_1 B_2 + A_2 B_1 + A_2 B_2 \\[3pt]
	V = (A_1 - A_2) \cdot (B_1 - B_2) = A_1 B_1 - A_1 B_2 - A_2 B_1 + A_2 B_2
\end{array} \qquad \begin{array}{l}
	C_1 = (U + V)/2 = A_1 B_1 + A_2 B_2 \\[3pt]
	C_2 = (U - V)/2 = A_1 B_2 + A_2 B_1
\end{array}\]
L'equazione alle ricorrenze per questo approccio è la seguente:
\[T(n) \; = \; \begin{cases}
	\; 1 & \text{se } n = 1 \\[3pt]
	\; 2 \cdot T(n/2) + 6 \cdot n/2 + 2 \cdot n/2 & \text{se } n > 1
\end{cases} \;\; = \;\; \begin{cases}
	\; 1 & \text{se } n = 1 \\[3pt]
	\; 2 \cdot T(n/2) + 4n & \text{se } n > 1
\end{cases}\]
Dal master theorem con \(a = 2\), \(b = 2\) e \(\omega(n) = 4n\) si ottiene una funzione soglia di
\(\sigma(n) = n^{\log_2 2} = n\) paragonabile a \(\omega(n)\). La complessità dal secondo caso del
master theorem risulta essere \(T(n) = \Theta(n \cdot \log n)\).

\subsection{Trasformata di Hadamard e Fast Hadamard Transform (FHT)}
\subsubsection*{Autovalori alle foglie dell'albero della ricorsione}
L'albero della ricorsione per il prodotto di matrici supercircolanti con approccio D\&C ottimizzato è un albero binario
completo in cui ogni nodo figlio sinistro riceve le somme tra i blocchi \(A_1\) e \(A_2\) e tra i blocchi \(B_1\) e \(B_2\),
mentre ogni nodo figlio destro riceve le differenze tra i medesimi blocchi, affinché siano moltiplicate opportunamente.

Le foglie ricevono due matrici degeneri in scalati \(a\) e \(b\) che è possibile dimostrare essere proprio gli autovalori
delle due matrici supercircolanti \(A\) e \(B\) di partenza.

\subsubsection*{Trasformata e prodotto di Hadamard}
È possibile interpretare le operazioni di divide come una trasformazione lineare che, partendo da una matrice supercircolante
rappresentata con la sua prima riga, restituisce il vettore di autovalori della matrice supercircolante stessa. In questo modo
il prodotto di matrici supercircolanti può essere interpretato come un prodotto di autovalori, che risulta avere complessità
\(\Theta(n)\) con \(n\) pari alla taglia delle matrici supercircolanti. Una volta fatto ciò, le operazioni di combine effettuano
l'inverso della trasformazione lineare, restituendo la matrice supercircolante risultante dal prodotto.

La trasformazione lineare impiegata prende il nome di \textbf{trasformata di Hadamard}, mentre il prodotto membro a membro
tra i vettori di autovalori prende il nome di \textbf{prodotto di Hadamard}.

La trasformata e la relativa antitrasformata di Hadamard possono essere calcolate in \(\Theta(n \cdot \log n)\), mentre
il prodotto di Hadamard può essere calcolato in \(\Theta(n)\). Di conseguenza, il prodotto di matrici supercircolanti tramite
la trasformata di Hadamard ha complessità complessiva di \(\Theta(n \cdot \log n)\).

\subsubsection*{Cambio di rappresentazioni per risolvere problemi più efficientemente}
L'idea alla base del prodotto di matrici supercircolanti tramite la trasformata di Hadamard è quello di convertire i dati
in ingresso in una rappresentazione che permette di risolvere il problema richiesto in maniera più efficiente, per poi
convertire il risultato ottenuto nella rappresentazione originale. Altri esempi di algoritmi che adottano questa tecnica sono:
\begin{itemize}
	\item l'algoritmo di Strassen, anche se non è chiaro come interpretare la rappresentazione intermedia
	\item l'algoritmo di FFT, che utilizza una trasformata di Fourier per calcolare il prodotto di polinomi in maniera
	più efficiente
\end{itemize}

\subsubsection*{Osservazione sulle simmetrie}
Una simmetria è una proprietà di invarianza rispetto ad una certa trasformazione. Ad esempio la proprietà commutativa è
l'invarianza rispetto alla trasformazione che scambia i due operandi di un'operazione. 

Conoscere le simmetrie permette di identificare ed eliminare i gradi di libertà di un problema, riducendo sia il peso della
rappresentazione utilizzata che nelle operazioni di calcolo.

\newpage

\section{Moltiplicazione di polinomi e FFT - Fast Fourier Transform}
